% Source: book pp. 59--68 (PDF 73--82); solutions pp. 65--82, 170--171 (A3).
\section{Null Spaces and Ranges}

\subsection{Null Space and Injectivity}

In this section we will learn about two subspaces that are intimately
connected with each linear map. We begin with the set of vectors that get
mapped to $0$.

\begin{definition}[Null space, $\nul T$]\label{ax:3.11}
For $T\in\Lin(V,W)$, the \emph{null space} of $T$, denoted by $\nul T$, is
the subset of $V$ consisting of those vectors that $T$ maps to $0$:
\[ \nul T=\{v\in V : Tv=0\}. \]
\end{definition}

\marginnote[14mm]{The word ``null'' means zero. Thus the term ``null space''
should remind you of the connection to $0$. Some mathematicians use the term
\emph{kernel} instead of null space.}%
\begin{example}[Null space]\label{ax:3.12}\leavevmode
\begin{itemize}
\item If $T$ is the zero map from $V$ to $W$, meaning that $Tv=0$ for every
  $v\in V$, then $\nul T=V$.
\item Suppose $\varphi\in\Lin(\C^3,\C)$ is defined by
  $\varphi(z_1,z_2,z_3)=z_1+2z_2+3z_3$. Then $\nul\varphi$ equals
  $\{(z_1,z_2,z_3)\in\C^3 : z_1+2z_2+3z_3=0\}$, which is a subspace of the
  domain of $\varphi$. We will soon see that the null space of each linear map
  is a subspace of its domain.
\item Suppose $D\in\Lin(\Poly(\R))$ is the differentiation map defined by
  $Dp=p'$. The only functions whose derivative equals the zero function are
  the constant functions. Thus the null space of $D$ equals the set of
  constant functions.
\item Suppose that $T\in\Lin(\Poly(\R))$ is the multiplication by $x^2$ map
  defined by $(Tp)(x)=x^2p(x)$. The only polynomial $p$ such that
  $x^2p(x)=0$ for all $x\in\R$ is the $0$ polynomial. Thus $\nul T=\{0\}$.
\item Suppose $T\in\Lin(\F^\infty)$ is the backward shift defined by
  \[ T(x_1,x_2,x_3,\dots)=(x_2,x_3,\dots). \]
  Then $T(x_1,x_2,x_3,\dots)$ equals $0$ if and only if the numbers
  $x_2,x_3,\dots$ are all $0$. Thus $\nul T=\{(a,0,0,\dots) : a\in\F\}$.
\end{itemize}
\end{example}

The next result shows that the null space of each linear map is a subspace of
the domain. In particular, $0$ is in the null space of every linear map.

\begin{result}[The null space is a subspace]\label{ax:3.13}
Suppose $T\in\Lin(V,W)$. Then $\nul T$ is a subspace of $V$.
\end{result}

\begin{proof}
Because $T$ is a linear map, $T(0)=0$ (by \ref{ax:3.10}). Thus
$0\in\nul T$.

Suppose $u,v\in\nul T$. Then
\[ T(u+v)=Tu+Tv=0+0=0. \]
Hence $u+v\in\nul T$. Thus $\nul T$ is closed under addition.

Suppose $u\in\nul T$ and $\lambda\in\F$. Then
\[ T(\lambda u)=\lambda Tu=\lambda0=0. \]
Hence $\lambda u\in\nul T$. Thus $\nul T$ is closed under scalar
multiplication.

We have shown that $\nul T$ contains $0$ and is closed under addition and
scalar multiplication. Thus $\nul T$ is a subspace of $V$ (by
\ref{ax:1.34}).
\end{proof}

As we will soon see, for a linear map the next definition is closely
connected to the null space.

\begin{definition}[Injective]\label{ax:3.14}
A function $T\colon V\to W$ is called \emph{injective} if $Tu=Tv$ implies
$u=v$.
\end{definition}

\marginnote{The term \emph{one-to-one} means the same as injective.}%
We could rephrase the definition above to say that $T$ is injective if
$u\ne v$ implies that $Tu\ne Tv$. Thus $T$ is injective if and only if it maps
distinct inputs to distinct outputs.

The next result says that we can check whether a linear map is injective by
checking whether $0$ is the only vector that gets mapped to $0$. As a simple
application of this result, we see that of the linear maps whose null spaces
we computed in \ref{ax:3.12}, only multiplication by $x^2$ is injective
(except that the zero map is injective in the special case $V=\{0\}$).

\begin{result}[Injectivity $\iff$ null space equals $\{0\}$]\label{ax:3.15}
Let $T\in\Lin(V,W)$. Then $T$ is injective if and only if $\nul T=\{0\}$.
\end{result}

\begin{proof}
First suppose $T$ is injective. We want to prove that $\nul T=\{0\}$. We
already know that $\{0\}\subseteq\nul T$ (by \ref{ax:3.10}). To prove the
inclusion in the other direction, suppose $v\in\nul T$. Then
\[ T(v)=0=T(0). \]
Because $T$ is injective, the equation above implies that $v=0$. Thus we can
conclude that $\nul T=\{0\}$, as desired.

To prove the implication in the other direction, now suppose
$\nul T=\{0\}$. We want to prove that $T$ is injective. To do this, suppose
$u,v\in V$ and $Tu=Tv$. Then
\[ 0=Tu-Tv=T(u-v). \]
Thus $u-v$ is in $\nul T$, which equals $\{0\}$. Hence $u-v=0$, which
implies that $u=v$. Hence $T$ is injective, as desired.
\end{proof}

\subsection{Range and Surjectivity}

Now we give a name to the set of outputs of a linear map.

\begin{definition}[Range]\label{ax:3.16}
For $T\in\Lin(V,W)$, the \emph{range} of $T$ is the subset of $W$ consisting
of those vectors that are equal to $Tv$ for some $v\in V$:
\[ \range T=\{Tv : v\in V\}. \]
\end{definition}

\begin{example}[Range]\label{ax:3.17}\leavevmode
\begin{itemize}
\item If $T$ is the zero map from $V$ to $W$, meaning that $Tv=0$ for every
  $v\in V$, then $\range T=\{0\}$.
\item Suppose $T\in\Lin(\R^2,\R^3)$ is defined by $T(x,y)=(2x,5y,x+y)$. Then
  \[ \range T=\bigl\{(2x,5y,x+y) : x,y\in\R\bigr\}. \]
  Note that $\range T$ is a subspace of $\R^3$. We will soon see that the
  range of each element of $\Lin(V,W)$ is a subspace of $W$.
\item Suppose $D\in\Lin(\Poly(\R))$ is the differentiation map defined by
  $Dp=p'$. Because for every polynomial $q\in\Poly(\R)$ there exists a
  polynomial $p\in\Poly(\R)$ such that $p'=q$, the range of $D$ is
  $\Poly(\R)$.
\end{itemize}
\end{example}

The next result shows that the range of each linear map is a subspace of the
vector space into which it is being mapped.

\begin{result}[The range is a subspace]\label{ax:3.18}
If $T\in\Lin(V,W)$, then $\range T$ is a subspace of $W$.
\end{result}

\begin{proof}
Suppose $T\in\Lin(V,W)$. Then $T(0)=0$ (by \ref{ax:3.10}), which implies that
$0\in\range T$.

If $w_1,w_2\in\range T$, then there exist $v_1,v_2\in V$ such that $Tv_1=w_1$
and $Tv_2=w_2$. Thus
\[ T(v_1+v_2)=Tv_1+Tv_2=w_1+w_2. \]
Hence $w_1+w_2\in\range T$. Thus $\range T$ is closed under addition.

If $w\in\range T$ and $\lambda\in\F$, then there exists $v\in V$ such that
$Tv=w$. Thus
\[ T(\lambda v)=\lambda Tv=\lambda w. \]
Hence $\lambda w\in\range T$. Thus $\range T$ is closed under scalar
multiplication.

We have shown that $\range T$ contains $0$ and is closed under addition and
scalar multiplication. Thus $\range T$ is a subspace of $W$ (by
\ref{ax:1.34}).
\end{proof}

\begin{definition}[Surjective]\label{ax:3.19}
A function $T\colon V\to W$ is called \emph{surjective} if its range equals
$W$.
\end{definition}

To illustrate the definition above, note that of the ranges we computed in
\ref{ax:3.17}, only the differentiation map is surjective (except that the
zero map is surjective in the special case $W=\{0\}$).

\marginnote{Some people use the term \emph{onto}, which means the same as
surjective.}%
Whether a linear map is surjective depends on what we are thinking of as the
vector space into which it maps.

\begin{example}[Surjectivity depends on the target space]\label{ax:3.20}
The differentiation map $D\in\Lin(\Poly_5(\R))$ defined by $Dp=p'$ is not
surjective, because the polynomial $x^5$ is not in the range of $D$. However,
the differentiation map $S\in\Lin(\Poly_5(\R),\Poly_4(\R))$ defined by
$Sp=p'$ is surjective, because its range equals $\Poly_4(\R)$, which is the
vector space into which $S$ maps.
\end{example}

\subsection{Fundamental Theorem of Linear Maps}

The next result is so important that it gets a dramatic name.

\begin{result}[Fundamental theorem of linear maps]\label{ax:3.21}
Suppose $V$ is finite-dimensional and $T\in\Lin(V,W)$. Then $\range T$ is
finite-dimensional and
\[ \dim V=\dim\nul T+\dim\range T. \]
\end{result}

\begin{proof}
Let $u_1,\dots,u_m$ be a basis of $\nul T$; thus $\dim\nul T=m$. The linearly
independent list $u_1,\dots,u_m$ can be extended to a basis
\[ u_1,\dots,u_m,v_1,\dots,v_n \]
of $V$ (by \ref{ax:2.32}). Thus $\dim V=m+n$. To complete the proof, we need to
show that $\range T$ is finite-dimensional and $\dim\range T=n$. We will do
this by proving that $Tv_1,\dots,Tv_n$ is a basis of $\range T$.

Let $v\in V$. Because $u_1,\dots,u_m,v_1,\dots,v_n$ spans $V$, we can write
\[ v=a_1u_1+\cdots+a_mu_m+b_1v_1+\cdots+b_nv_n, \]
where the $a$'s and $b$'s are in $\F$. Applying $T$ to both sides of this
equation, we get
\[ Tv=b_1Tv_1+\cdots+b_nTv_n, \]
where the terms of the form $Tu_k$ disappeared because each $u_k$ is in
$\nul T$. The last equation implies that the list $Tv_1,\dots,Tv_n$ spans
$\range T$. In particular, $\range T$ is finite-dimensional.

To show $Tv_1,\dots,Tv_n$ is linearly independent, suppose
$c_1,\dots,c_n\in\F$ and
\[ c_1Tv_1+\cdots+c_nTv_n=0. \]
Then
\[ T(c_1v_1+\cdots+c_nv_n)=0. \]
Hence
\[ c_1v_1+\cdots+c_nv_n\in\nul T. \]
Because $u_1,\dots,u_m$ spans $\nul T$, we can write
\[ c_1v_1+\cdots+c_nv_n=d_1u_1+\cdots+d_mu_m, \]
where the $d$'s are in $\F$. This equation implies that all the $c$'s (and
$d$'s) are $0$ (because $u_1,\dots,u_m,v_1,\dots,v_n$ is linearly
independent). Thus $Tv_1,\dots,Tv_n$ is linearly independent and hence is a
basis of $\range T$, as desired.
\end{proof}

Now we can show that no linear map from a finite-dimensional vector space to a
``smaller'' vector space can be injective, where ``smaller'' is measured by
dimension.

\begin{result}[Linear map to a lower-dimensional space is not injective]\label{ax:3.22}
Suppose $V$ and $W$ are finite-dimensional vector spaces such that
$\dim V>\dim W$. Then no linear map from $V$ to $W$ is injective.
\end{result}

\begin{proof}
Let $T\in\Lin(V,W)$. Then
\begin{align*}
\dim\nul T &= \dim V-\dim\range T\\
  &\ge \dim V-\dim W\\
  &> 0,
\end{align*}
where the first line above comes from the fundamental theorem of linear maps
(\ref{ax:3.21}) and the second line follows from \ref{ax:2.37}. The inequality
above states that $\dim\nul T>0$. This means that $\nul T$ contains vectors
other than $0$. Thus $T$ is not injective (by \ref{ax:3.15}).
\end{proof}

\begin{example}[Linear map from $\F^4$ to $\F^3$ is not injective]\label{ax:3.23}
Define a linear map $T\colon\F^4\to\F^3$ by
\[ T(z_1,z_2,z_3,z_4)=\bigl(\sqrt7z_1+\pi z_2+z_4,\,97z_1+3z_2+2z_3,\,
   z_2+6z_3+7z_4\bigr). \]
Because $\dim\F^4>\dim\F^3$, we can use \ref{ax:3.22} to assert that $T$ is
not injective, without doing any calculations.
\end{example}

The next result shows that no linear map from a finite-dimensional vector
space to a ``bigger'' vector space can be surjective, where ``bigger'' is
measured by dimension.

\begin{result}[Linear map to a higher-dimensional space is not surjective]\label{ax:3.24}
Suppose $V$ and $W$ are finite-dimensional vector spaces such that
$\dim V<\dim W$. Then no linear map from $V$ to $W$ is surjective.
\end{result}

\begin{proof}
Let $T\in\Lin(V,W)$. Then
\begin{align*}
\dim\range T &= \dim V-\dim\nul T\\
  &\le \dim V\\
  &< \dim W,
\end{align*}
where the equality above comes from the fundamental theorem of linear maps
(\ref{ax:3.21}). The inequality above states that $\dim\range T<\dim W$. This
means that $\range T$ cannot equal $W$. Thus $T$ is not surjective.
\end{proof}

As we will soon see, \ref{ax:3.22} and \ref{ax:3.24} have important
consequences in the theory of linear equations. The idea is to express
questions about systems of linear equations in terms of linear maps. Let's
begin by rephrasing in terms of linear maps the question of whether a
homogeneous system of linear equations has a nonzero solution.

\marginnote{\emph{Homogeneous}, in this context, means that the constant term
on the right side of each equation below is $0$.}%
Fix positive integers $m$ and $n$, and let $A_{j,k}\in\F$ for $j=1,\dots,m$
and $k=1,\dots,n$. Consider the homogeneous system of linear equations
\[ \begin{gathered}
   \sum_{k=1}^n A_{1,k}x_k=0\\
   \vdots\\
   \sum_{k=1}^n A_{m,k}x_k=0.
   \end{gathered} \]
Clearly $x_1=\cdots=x_n=0$ is a solution of the system of equations above;
the question here is whether any other solutions exist.

Define $T\colon\F^n\to\F^m$ by
\[ \refstepcounter{theorem}\tag*{\thetheorem}\label{ax:3.25}
   T(x_1,\dots,x_n)=\Bigl(\sum_{k=1}^n A_{1,k}x_k,\dots,
   \sum_{k=1}^n A_{m,k}x_k\Bigr). \]
The equation $T(x_1,\dots,x_n)=0$ (the $0$ here is the additive identity in
$\F^m$, namely, the list of length $m$ of all $0$'s) is the same as the
homogeneous system of linear equations above.

Thus we want to know if $\nul T$ is strictly bigger than $\{0\}$, which is
equivalent to $T$ not being injective (by \ref{ax:3.15}). The next result
gives an important condition for ensuring that $T$ is not injective.

\begin{result}[Homogeneous system of linear equations]\label{ax:3.26}
A homogeneous system of linear equations with more variables than equations
has nonzero solutions.
\end{result}

\begin{proof}
Use the notation and result from the discussion above. Thus $T$ is a linear
map from $\F^n$ to $\F^m$, and we have a homogeneous system of $m$ linear
equations with $n$ variables $x_1,\dots,x_n$. From \ref{ax:3.22} we see that
$T$ is not injective if $n>m$.
\end{proof}

Example of the result above: a homogeneous system of four linear equations
with five variables has nonzero solutions.

Now we consider the question of whether a system of linear equations has no
solutions for some choice of the constant terms. To rephrase this question in
terms of a linear map, fix positive integers $m$ and $n$, and let
$A_{j,k}\in\F$ for all $j=1,\dots,m$ and all $k=1,\dots,n$. For
$c_1,\dots,c_m\in\F$, consider the system of linear equations
\[ \refstepcounter{theorem}\tag*{\thetheorem}\label{ax:3.27}
   \begin{gathered}
   \sum_{k=1}^n A_{1,k}x_k=c_1\\
   \vdots\\
   \sum_{k=1}^n A_{m,k}x_k=c_m.
   \end{gathered} \]
With this notation, the question here is whether there is some choice of the
constant terms $c_1,\dots,c_m\in\F$ such that no solution exists to the system
above.

\marginnote{The results \ref{ax:3.26} and \ref{ax:3.28}, which compare the
number of variables and the number of equations, can also be proved using
Gaussian elimination. The abstract approach taken here seems to provide
cleaner proofs.}%
Define $T\colon\F^n\to\F^m$ as in \ref{ax:3.25}. The equation
$T(x_1,\dots,x_n)=(c_1,\dots,c_m)$ is the same as the system of equations
\ref{ax:3.27}. Thus we want to know if $\range T\ne\F^m$. Hence we can
rephrase our question about not having a solution for some choice of
$c_1,\dots,c_m\in\F$ as follows: What condition ensures that $T$ is not
surjective? The next result gives one such condition.

\begin{result}[System of linear equations with more equations than variables]\label{ax:3.28}
A system of linear equations with more equations than variables has no
solution for some choice of the constant terms.
\end{result}

\begin{proof}
Use the notation from the discussion above. Thus $T$ is a linear map from
$\F^n$ to $\F^m$, and we have a system of $m$ equations with $n$ variables
$x_1,\dots,x_n$; see \ref{ax:3.27}. If $n<m$, then \ref{ax:3.24} implies that
$T$ is not surjective. As discussed above, this shows that if we have more
equations than variables in a system of linear equations, then there is no
solution for some choice of the constant terms.
\end{proof}

Example of the result above: a system of five linear equations with four
variables has no solution for some choice of the constant terms.

% ------------------------------------------------------------------ exercises
\exercisesheading{3B}

\begin{exercise}
Give an example of a linear map $T$ with $\dim\nul T=3$ and
$\dim\range T=2$.
\end{exercise}
\begin{solution}
Let $T\in\Lin(\R^5,\R^2)$ be defined by
\[ T(x_1,x_2,x_3,x_4,x_5)=(x_1,x_2). \]
Then $\nul T=\{(0,0,x_3,x_4,x_5)\mid x_3,x_4,x_5\in\R\}$, which has dimension
$3$, and $\range T=\R^2$, which has dimension $2$.
\end{solution}

\begin{exercise}
Suppose $S,T\in\Lin(V)$ are such that $\range S\subseteq\nul T$. Prove that
$(ST)^2=0$.
\end{exercise}
\begin{solution}
Let $v\in V$. Then
\[ (ST)^2v=ST(STv)=S\bigl(T(STv)\bigr). \]
Since $STv\in\range S$ and $\range S\subseteq\nul T$, we have $T(STv)=0$.
Therefore,
\[ (ST)^2v=S(0)=0. \]
Since $v$ was arbitrary, it follows that $(ST)^2=0$.
\end{solution}

\begin{exercise}
Suppose $v_1,\dots,v_m$ is a list of vectors in $V$. Define
$T\in\Lin(\F^m,V)$ by
\[ T(z_1,\dots,z_m)=z_1v_1+\cdots+z_mv_m. \]
\begin{parts}
\item What property of $T$ corresponds to $v_1,\dots,v_m$ spanning $V$?
\item What property of $T$ corresponds to the list $v_1,\dots,v_m$ being
  linearly independent?
\end{parts}
\end{exercise}
\begin{solution}\leavevmode
\begin{parts}
\item The list $v_1,\dots,v_m$ spans $V$ if and only if $\range T=V$, i.e.,
  $T$ is surjective.
\item The list $v_1,\dots,v_m$ is linearly independent if and only if
  $\nul T=\{0\}$, i.e., $T$ is injective.
\end{parts}
\end{solution}

\begin{exercise}
Show that $\{T\in\Lin(\R^5,\R^4) : \dim\nul T>2\}$ is not a subspace of
$\Lin(\R^5,\R^4)$.
\end{exercise}
\begin{solution}
Let $T_1,T_2\in\Lin(\R^5,\R^4)$ be defined by
\[ T_1(x_1,x_2,x_3,x_4,x_5)=(x_1,x_2,0,0) \]
and
\[ T_2(x_1,x_2,x_3,x_4,x_5)=(0,0,x_3,x_4). \]
Then $\nul T_1=\{(0,0,x_3,x_4,x_5)\mid x_3,x_4,x_5\in\R\}$ and
$\nul T_2=\{(x_1,x_2,0,0,x_5)\mid x_1,x_2,x_5\in\R\}$ both have dimension
$3$, so $T_1$ and $T_2$ are in the set
$\{T\in\Lin(\R^5,\R^4)\mid\dim\nul T>2\}$. However,
\[ (T_1+T_2)(x_1,x_2,x_3,x_4,x_5)=(x_1,x_2,x_3,x_4), \]
which has null space $\nul(T_1+T_2)=\{(0,0,0,0,x_5)\mid x_5\in\R\}$ of
dimension $1$. Therefore,
\[ \dim\nul(T_1+T_2)=1<2, \]
so $T_1+T_2$ is not in the set. Hence, the set is not closed under addition
and is not a subspace of $\Lin(\R^5,\R^4)$.
\end{solution}

\begin{exercise}
Give an example of $T\in\Lin(\R^4)$ such that $\range T=\nul T$.
\end{exercise}
\begin{solution}
Let $v_1,v_2,v_3,v_4$ be the standard basis of $\R^4$. Define
$T\in\Lin(\R^4)$ by
\[ Tv_1=v_3, \quad Tv_2=v_4, \quad Tv_3=0, \quad Tv_4=0. \]
Then for any $x=(x_1,x_2,x_3,x_4)\in\R^4$,
\[ Tx=x_1Tv_1+x_2Tv_2+x_3Tv_3+x_4Tv_4=x_1v_3+x_2v_4. \]
Then $\range T=\spn v_3,v_4$. If $Tx=0$, then $x_1v_3+x_2v_4=0$, which
implies $x_1=x_2=0$. Therefore,
$\nul T=\{(0,0,x_3,x_4)\mid x_3,x_4\in\R\}=\spn v_3,v_4$. Hence,
$\range T=\nul T$.
\end{solution}

\begin{exercise}
Prove that there does not exist $T\in\Lin(\R^5)$ such that
$\range T=\nul T$.
\end{exercise}
\begin{solution}
Suppose for contradiction that there exists $T\in\Lin(\R^5)$ such that
$\range T=\nul T$. Then by the fundamental theorem of linear maps,
\[ \dim\R^5=\dim\nul T+\dim\range T=2\dim\range T. \]
Therefore, $5=2\dim\range T$, which is impossible since $5$ is odd and
$2\dim\range T$ is even. Hence, no such $T$ exists.
\end{solution}

\begin{exercise}
Suppose $V$ and $W$ are finite-dimensional with $2\le\dim V\le\dim W$. Show
that $\{T\in\Lin(V,W) : T \text{ is not injective}\}$ is not a subspace of
$\Lin(V,W)$.
\end{exercise}
\begin{solution}
Let $\dim V=n$ and $\dim W=m$ with $2\le n\le m$. Let $v_1,v_2,\dots,v_n$ be a
basis of $V$, and let $w_1,w_2,\dots,w_m$ be a basis of $W$. Define
$T_1,T_2\in\Lin(V,W)$ by
\[ T_1v_1=w_1, \quad T_1v_k=0 \text{ for } k\ge2, \]
and
\[ T_2v_1=0, \quad T_2v_k=w_k \text{ for } k\ge2. \]
Then $v_2\in\nul T_1$ and $v_1\in\nul T_2$, so both $T_1$ and $T_2$ are not
injective. However,
\[ (T_1+T_2)v_1=T_1v_1+T_2v_1=w_1+0=w_1\ne0, \]
and for $k\ge2$,
\[ (T_1+T_2)v_k=T_1v_k+T_2v_k=0+w_k=w_k\ne0. \]
Suppose $v\in V$ is such that $(T_1+T_2)v=0$. Then $v$ can be expressed as a
linear combination of the basis vectors:
\[ v=\alpha_1v_1+\alpha_2v_2+\cdots+\alpha_nv_n. \]
Then
\[ (T_1+T_2)v=\alpha_1w_1+\alpha_2w_2+\cdots+\alpha_nw_n=0. \]
Since $w_1,w_2,\dots,w_m$ are linearly independent, it follows that
$\alpha_1=\alpha_2=\cdots=\alpha_n=0$, hence $v=0$. Therefore, $T_1+T_2$ is
injective. Since $T_1$ and $T_2$ are not injective, but their sum is
injective, the set $\{T\in\Lin(V,W)\mid T \text{ is not injective}\}$ is not
closed under addition and thus is not a subspace of $\Lin(V,W)$.
\end{solution}

\begin{exercise}
Suppose $V$ and $W$ are finite-dimensional with $\dim V\ge\dim W\ge2$. Show
that $\{T\in\Lin(V,W) : T \text{ is not surjective}\}$ is not a subspace of
$\Lin(V,W)$.
\end{exercise}
\begin{solution}
We proceed similarly to the previous exercise. Let $\dim V=n$ and $\dim W=m$
with $n\ge m\ge2$. Let $v_1,v_2,\dots,v_n$ be a basis of $V$, and let
$w_1,w_2,\dots,w_m$ be a basis of $W$. Define $T_1,T_2\in\Lin(V,W)$ by
\[ T_1v_1=w_1, \quad T_1v_k=0 \text{ for } k\ge2, \]
and
\[ T_2v_1=0, \quad T_2v_k=w_k \text{ for } k=2,\dots,m, \quad
   T_2v_k=0 \text{ for } k>m. \]
Then $\range T_1=\spn w_1$ and $\range T_2=\spn w_2,\dots,w_m$, both of which
are proper subspaces of $W$, so $T_1$ and $T_2$ are not surjective. However,
\[ (T_1+T_2)v_1=T_1v_1+T_2v_1=w_1+0=w_1, \]
and for $k=2,\dots,m$,
\[ (T_1+T_2)v_k=T_1v_k+T_2v_k=0+w_k=w_k. \]
Therefore, $\range(T_1+T_2)=\spn w_1,w_2,\dots,w_m=W$, so $T_1+T_2$ is
surjective. Since $T_1$ and $T_2$ are not surjective, but their sum is
surjective, the set $\{T\in\Lin(V,W)\mid T \text{ is not surjective}\}$ is not
closed under addition and thus is not a subspace of $\Lin(V,W)$.
\end{solution}

\begin{exercise}
Suppose $T\in\Lin(V,W)$ is injective and $v_1,\dots,v_n$ is linearly
independent in $V$. Prove that $Tv_1,\dots,Tv_n$ is linearly independent in
$W$.
\end{exercise}
\begin{solution}
Suppose $a_1,\dots,a_n\in\F$ are such that
\[ a_1Tv_1+\cdots+a_nTv_n=0. \]
By linearity of $T$, we have
\[ T(a_1v_1+\cdots+a_nv_n)=0. \]
Injectivity of $T$ implies that $\nul T=\{0\}$, so
\[ a_1v_1+\cdots+a_nv_n=0. \]
Since $v_1,\dots,v_n$ is linearly independent, it follows that
$a_1=\cdots=a_n=0$. Therefore, $Tv_1,\dots,Tv_n$ is linearly independent in
$W$.
\end{solution}

\begin{exercise}
Suppose $v_1,\dots,v_n$ spans $V$ and $T\in\Lin(V,W)$. Show that
$Tv_1,\dots,Tv_n$ spans $\range T$.
\end{exercise}
\begin{solution}
Let $w\in\range T$. Then there exists $v\in V$ such that $Tv=w$. Since
$v_1,\dots,v_n$ spans $V$, we can write
\[ v=a_1v_1+\cdots+a_nv_n \]
for some scalars $a_1,\dots,a_n\in\F$. By linearity of $T$, we have
\[ w=Tv=T(a_1v_1+\cdots+a_nv_n)=a_1Tv_1+\cdots+a_nTv_n. \]
Therefore, $w$ is a linear combination of $Tv_1,\dots,Tv_n$. Since $w$ was
arbitrary in $\range T$, it follows that $Tv_1,\dots,Tv_n$ spans $\range T$.
\end{solution}

\begin{exercise}
Suppose that $V$ is finite-dimensional and that $T\in\Lin(V,W)$. Prove that
there exists a subspace $U$ of $V$ such that
\[ U\cap\nul T=\{0\} \quad\text{and}\quad \range T=\{Tu : u\in U\}. \]
\end{exercise}
\begin{solution}
\solnote{Manual Remark}{Note that this exercise proves the existence of a
subspace $U$ of $V$ such that $V=U\oplus\nul T$.}

Since $\nul T$ is a subspace of $V$, we can extend a basis of $\nul T$ to a
basis of $V$. Let $n_1,\dots,n_k$ be a basis of $\nul T$, and extend this to
a basis of $V$ by adding vectors $u_1,\dots,u_m$. Hence, a basis for $V$ is
given by
\[ n_1,\dots,n_k,u_1,\dots,u_m. \]
Define $U=\spn(u_1,\dots,u_m)$. Suppose $v\in U\cap\nul T$. Then $v$ can be
expressed as a linear combination of the basis vectors of $U$:
\[ v=\alpha_1u_1+\cdots+\alpha_mu_m. \]
Since $v\in\nul T$ as well, we can also express $v$ as a linear combination
of the basis vectors of $\nul T$:
\[ v=\beta_1n_1+\cdots+\beta_kn_k. \]
Equating the two expressions for $v$, we have
\[ \alpha_1u_1+\cdots+\alpha_mu_m=\beta_1n_1+\cdots+\beta_kn_k. \]
Since $n_1,\dots,n_k,u_1,\dots,u_m$ is a basis for $V$, the only way this
equality can hold is if all coefficients are zero, i.e.,
$\alpha_1=\cdots=\alpha_m=0$ and $\beta_1=\cdots=\beta_k=0$. Therefore,
$v=0$, and we conclude that $U\cap\nul T=\{0\}$.

Now, let $w\in\range T$. Then there exists $v\in V$ such that $Tv=w$. We can
express $v$ as a linear combination of the basis vectors of $V$:
\[ v=\gamma_1n_1+\cdots+\gamma_kn_k+\delta_1u_1+\cdots+\delta_mu_m. \]
Applying $T$ to both sides, we have
\begin{align*}
Tv &= T(\gamma_1n_1+\cdots+\gamma_kn_k+\delta_1u_1+\cdots+\delta_mu_m)\\
   &= \gamma_1Tn_1+\cdots+\gamma_kTn_k+\delta_1Tu_1+\cdots+\delta_mTu_m.
\end{align*}
Since $n_1,\dots,n_k\in\nul T$, we have $Tn_i=0$ for all $i=1,\dots,k$.
Therefore,
\[ Tv=\delta_1Tu_1+\cdots+\delta_mTu_m. \]
Let $u=\delta_1u_1+\cdots+\delta_mu_m\in U$. Then $Tu=w$. Since $w$ was
arbitrary in $\range T$, it follows that $\range T=\{Tu\mid u\in U\}$. Thus,
we have found a subspace $U$ of $V$ such that $U\cap\nul T=\{0\}$ and
$\range T=\{Tu\mid u\in U\}$.
\end{solution}

\begin{exercise}
Suppose $T$ is a linear map from $\F^4$ to $\F^2$ such that
\[ \nul T=\{(x_1,x_2,x_3,x_4)\in\F^4 : x_1=5x_2 \text{ and } x_3=7x_4\}. \]
Prove that $T$ is surjective.
\end{exercise}
\begin{solution}
We may write the null space of $T$ as
\[ \nul T=\{(5x_2,x_2,7x_4,x_4)\mid x_2,x_4\in\F\}. \]
From this, we can see that $\nul T$ is spanned by the vectors $(5,1,0,0)$ and
$(0,0,7,1)$. Therefore, $\dim\nul T=2$. By the fundamental theorem of linear
maps, we have
\[ \dim\F^4=\dim\nul T+\dim\range T\implies4=2+\dim\range T. \]
Solving for $\dim\range T$, we find that $\dim\range T=2$. Since
$\dim\F^2=2$, it follows that $\range T=\F^2$. Therefore, $T$ is surjective.
\end{solution}

\begin{exercise}
Suppose $U$ is a three-dimensional subspace of $\R^8$ and that $T$ is a
linear map from $\R^8$ to $\R^5$ such that $\nul T=U$. Prove that $T$ is
surjective.
\end{exercise}
\begin{solution}
By the fundamental theorem of linear maps, we have
\[ \dim\R^8=\dim\nul T+\dim\range T\implies8=3+\dim\range T. \]
Solving for $\dim\range T$, we find that $\dim\range T=5$. Since
$\dim\R^5=5$, it follows that $\range T=\R^5$. Therefore, $T$ is surjective.
\end{solution}

\begin{exercise}
Prove that there does not exist a linear map from $\F^5$ to $\F^2$ whose null
space equals
$\{(x_1,x_2,x_3,x_4,x_5)\in\F^5 : x_1=3x_2 \text{ and } x_3=x_4=x_5\}$.
\end{exercise}
\begin{solution}
We can express the null space as
\[ \nul T=\{(3x_2,x_2,x_4,x_4,x_4)\mid x_2,x_4\in\F\}. \]
From this, we can see that $\nul T$ is spanned by the vectors $(3,1,0,0,0)$
and $(0,0,1,1,1)$. Therefore, $\dim\nul T=2$. By the fundamental theorem of
linear maps, we have
\[ \dim\F^5=\dim\nul T+\dim\range T\implies5=2+\dim\range T. \]
Solving for $\dim\range T$, we find that $\dim\range T=3$. However,
$\dim\F^2=2$, which is less than $3$. This is a contradiction, as the
dimension of the range cannot exceed the dimension of the target space.
Therefore, no such linear map exists.
\end{solution}

\begin{exercise}
Suppose there exists a linear map on $V$ whose null space and range are both
finite-dimensional. Prove that $V$ is finite-dimensional.
\end{exercise}
\begin{solution}
\solnote{Manual Remark}{The spanning list in the solution is actually a basis
of $V$.}

Let $T\in\Lin(V)$ be such that $\nul T$ and $\range T$ are both
finite-dimensional. Let $u_1,\dots,u_m$ be a basis of $\nul T$, and let
$w_1,\dots,w_n$ be a basis of $\range T$. Select $v_j\in V$ be such that
$Tv_j=w_j$ for $j=1,\dots,n$. We claim that the list
$u_1,\dots,u_m,v_1,\dots,v_n$ spans $V$.

Let $v\in V$. Then $Tv\in\range T$, so we may write
\[ Tv=\alpha_1w_1+\cdots+\alpha_nw_n \]
for some scalars $\alpha_1,\dots,\alpha_n\in\F$. Consider the vector
\[ u=v-(\alpha_1v_1+\cdots+\alpha_nv_n). \]
Then
\[ Tu=Tv-(\alpha_1Tv_1+\cdots+\alpha_nTv_n)=Tv-(\alpha_1w_1+\cdots+\alpha_nw_n)
   =0. \]
Therefore, $u\in\nul T$, and we can write
\[ u=\beta_1u_1+\cdots+\beta_mu_m \]
for some scalars $\beta_1,\dots,\beta_m\in\F$. Thus,
\[ v=u+(\alpha_1v_1+\cdots+\alpha_nv_n)
   =\beta_1u_1+\cdots+\beta_mu_m+\alpha_1v_1+\cdots+\alpha_nv_n, \]
which shows that the list spans $V$. Hence, $V$ is finite-dimensional.
\end{solution}

\begin{exercise}
Suppose $V$ and $W$ are both finite-dimensional. Prove that there exists an
injective linear map from $V$ to $W$ if and only if $\dim V\le\dim W$.
\end{exercise}
\begin{solution}
Suppose there exists an injective linear map $T\in\Lin(V,W)$. By the
fundamental theorem of linear maps, we have
\[ \dim V=\dim\nul T+\dim\range T. \]
Since $T$ is injective, $\nul T=\{0\}$, so $\dim\nul T=0$. Therefore,
\[ \dim V=0+\dim\range T=\dim\range T. \]
Since $\range T$ is a subspace of $W$, we have $\dim\range T\le\dim W$. Thus,
\[ \dim V=\dim\range T\le\dim W. \]

Conversely, suppose $\dim V\le\dim W$. Let $n=\dim V$ and $m=\dim W$. Choose
a basis $\{v_1,\dots,v_n\}$ of $V$. Let $\{w_1,\dots,w_m\}$ be a basis of $W$.
Define a linear map $T\in\Lin(V,W)$ by
\[ Tv_i=w_i \quad\text{for } i=1,\dots,n. \]
Since $n\le m$, we can assign each basis vector of $V$ to a distinct basis
vector of $W$. To show that $T$ is injective, let
$v=\alpha_1v_1+\cdots+\alpha_nv_n\in V$ such that $Tv=0$. Then
\[ Tv=\alpha_1Tv_1+\cdots+\alpha_nTv_n=\alpha_1w_1+\cdots+\alpha_nw_n=0. \]
Since $w_1,\dots,w_n$ are linearly independent, it follows that
$\alpha_1=\cdots=\alpha_n=0$. Therefore, $v=0$, and $T$ is injective. Hence,
there exists an injective linear map from $V$ to $W$ if and only if
$\dim V\le\dim W$.
\end{solution}

\begin{exercise}
Suppose $V$ and $W$ are both finite-dimensional. Prove that there exists a
surjective linear map from $V$ onto $W$ if and only if $\dim V\ge\dim W$.
\end{exercise}
\begin{solution}
Suppose there exists a surjective linear map $T\in\Lin(V,W)$. By the
fundamental theorem of linear maps, we have
\[ \dim V=\dim\nul T+\dim\range T. \]
Since $T$ is surjective, $\range T=W$, so $\dim\range T=\dim W$. Therefore,
\[ \dim V=\dim\nul T+\dim W. \]
Since $\dim\nul T\ge0$, it follows that
\[ \dim V\ge\dim W. \]

Conversely, suppose $\dim V\ge\dim W$. Let $n=\dim V$ and $m=\dim W$. Choose
a basis $\{v_1,\dots,v_n\}$ of $V$. Let $\{w_1,\dots,w_m\}$ be a basis of $W$.
Define a linear map $T\in\Lin(V,W)$ by
\[ Tv_i=w_i \quad\text{for } i=1,\dots,m, \qquad
   Tv_i=0 \quad\text{for } i=m+1,\dots,n. \]
To show that $T$ is surjective, let $w\in W$. Then $w$ can be expressed as a
linear combination of the basis vectors of $W$:
\[ w=\beta_1w_1+\cdots+\beta_mw_m. \]
Consider the vector $v=\beta_1v_1+\cdots+\beta_mv_m\in V$. Then
\[ Tv=\beta_1Tv_1+\cdots+\beta_mTv_m=\beta_1w_1+\cdots+\beta_mw_m=w. \]
Since $w$ was arbitrary in $W$, it follows that $\range T=W$. Therefore, $T$
is surjective. Hence, there exists a surjective linear map from $V$ onto $W$
if and only if $\dim V\ge\dim W$.
\end{solution}

\begin{exercise}
Suppose $V$ and $W$ are finite-dimensional and that $U$ is a subspace of $V$.
Prove that there exists $T\in\Lin(V,W)$ such that $\nul T=U$ if and only if
$\dim U\ge\dim V-\dim W$.
\end{exercise}
\begin{solution}
Suppose there exists $T\in\Lin(V,W)$ such that $\nul T=U$. By the
fundamental theorem of linear maps, we have
\[ \dim V=\dim\nul T+\dim\range T=\dim U+\dim\range T. \]
Since $\range T$ is a subspace of $W$, we have $\dim\range T\le\dim W$.
Therefore,
\[ \dim V=\dim U+\dim\range T\le\dim U+\dim W. \]
Rearranging this inequality gives
\[ \dim U\ge\dim V-\dim W. \]

Conversely, suppose $\dim U\ge\dim V-\dim W$. Let $n=\dim V$, $m=\dim W$, and
$k=\dim U$. Let $u_1,\dots,u_k$ be a basis of $U$. Extend this basis to a
basis of $V$ by adding vectors $v_1,\dots,v_{n-k}$. Define a linear map
$T\in\Lin(V,W)$ by
\[ Tu_i=0 \quad\text{for } i=1,\dots,k, \qquad
   Tv_j=w_j \quad\text{for } j=1,\dots,n-k, \]
where $w_1,\dots,w_{n-k}$ are linearly independent vectors in $W$. Since
$\dim U=k$ and $\dim V-\dim W=n-m$, we have $k\ge n-m$, which implies that
$n-k\le m$. Therefore, we can choose $w_1,\dots,w_{n-k}$ to be linearly
independent in $W$. To show that $\nul T=U$, let $v\in\nul T$. Then $Tv=0$.
Express $v$ as a linear combination of the basis vectors of $V$:
\[ v=\alpha_1u_1+\cdots+\alpha_ku_k+\beta_1v_1+\cdots+\beta_{n-k}v_{n-k}. \]
Applying $T$ to both sides, we have
\begin{align*}
Tv &= \alpha_1Tu_1+\cdots+\alpha_kTu_k+\beta_1Tv_1+\cdots+\beta_{n-k}Tv_{n-k}\\
   &= 0+\beta_1w_1+\cdots+\beta_{n-k}w_{n-k}=0.
\end{align*}
Since $w_1,\dots,w_{n-k}$ are linearly independent, it follows that
$\beta_1=\cdots=\beta_{n-k}=0$. Therefore,
\[ v=\alpha_1u_1+\cdots+\alpha_ku_k\in U. \]
Hence, $\nul T\subseteq U$. Conversely, if $v\in U$, then $Tv=0$ by the
definition of $T$. Therefore, $U\subseteq\nul T$. Combining these two
inclusions, we have $\nul T=U$. Thus, there exists $T\in\Lin(V,W)$ such that
$\nul T=U$ if and only if $\dim U\ge\dim V-\dim W$.
\end{solution}

\begin{exercise}
Suppose $W$ is finite-dimensional and $T\in\Lin(V,W)$. Prove that $T$ is
injective if and only if there exists $S\in\Lin(W,V)$ such that $ST$ is the
identity operator on $V$.
\end{exercise}
\begin{solution}
Suppose $T$ is injective. Then $\nul T=\{0\}$. Also, $\range T$ is finite
dimensional since it is a subspace of $W$. Since $T$ is injective, it is an
isomorphism from $V$ onto $\range T$. Hence, $V$ is finite dimensional. Let
$n=\dim V$ and $m=\dim W$. Since $T$ is injective, we have $\dim\range T=n$.
Choose a basis $\{v_1,\dots,v_n\}$ of $V$. Then $\{Tv_1,\dots,Tv_n\}$ is
linearly independent in $W$. Extend this set to a basis of $W$ by adding
vectors $w_1,\dots,w_{m-n}$. Define a linear map $S\in\Lin(W,V)$ by
\[ S(Tv_i)=v_i \quad\text{for } i=1,\dots,n, \qquad
   S(w_j)=0 \quad\text{for } j=1,\dots,m-n. \]
Then for any $v\in V$, we can express $v$ as a linear combination of the
basis vectors:
\[ v=\alpha_1v_1+\cdots+\alpha_nv_n. \]
Applying $T$ to both sides, we have
\[ Tv=\alpha_1Tv_1+\cdots+\alpha_nTv_n. \]
Now, applying $S$ to $Tv$, we get
\[ STv=S(\alpha_1Tv_1+\cdots+\alpha_nTv_n)=\alpha_1S(Tv_1)+\cdots
   +\alpha_nS(Tv_n)=\alpha_1v_1+\cdots+\alpha_nv_n=v. \]
Therefore, $ST$ is the identity operator on $V$.

Conversely, suppose there exists $S\in\Lin(W,V)$ such that $ST$ is the
identity operator on $V$. Let $v\in\nul T$. Then $Tv=0$, and applying $S$
gives
\[ STv=S(0)=0. \]
Since $ST$ is the identity operator on $V$, we have $v=STv=0$. Therefore,
$\nul T=\{0\}$, which means $T$ is injective.
\end{solution}

\begin{exercise}
Suppose $W$ is finite-dimensional and $T\in\Lin(V,W)$. Prove that $T$ is
surjective if and only if there exists $S\in\Lin(W,V)$ such that $TS$ is the
identity operator on $W$.
\end{exercise}
\begin{solution}
Suppose $T$ is surjective. Then $\range T=W$. Let $m=\dim W$ and choose a basis
$\{w_1,\dots,w_m\}$ of $W$. Since $T$ is surjective, for each $w_i$, there
exists $v_i\in V$ such that $Tv_i=w_i$. Define a linear map $S\in\Lin(W,V)$ by
\[ S(w_i)=v_i \quad\text{for } i=1,\dots,m. \]
Then for any $w\in W$, we can express $w$ as a linear combination of the basis
vectors:
\[ w=\alpha_1w_1+\cdots+\alpha_mw_m. \]
Applying $S$ to both sides, we have
\[ Sw=S(\alpha_1w_1+\cdots+\alpha_mw_m)=\alpha_1S(w_1)+\cdots+\alpha_mS(w_m)
   =\alpha_1v_1+\cdots+\alpha_mv_m. \]
Now, applying $T$ to $Sw$, we get
\[ TSw=T(\alpha_1v_1+\cdots+\alpha_mv_m)=\alpha_1Tv_1+\cdots+\alpha_mTv_m
   =\alpha_1w_1+\cdots+\alpha_mw_m=w. \]
Therefore, $TS$ is the identity operator on $W$.

Conversely, suppose there exists $S\in\Lin(W,V)$ such that $TS$ is the identity
operator on $W$. Let $w\in W$. Then
\[ TSw=w. \]
This means that for every $w\in W$, there exists a vector $v=Sw\in V$ such that
$Tv=T(Sw)=w$. Therefore, $\range T=W$, which means that $T$ is surjective.
\end{solution}

\begin{exercise}
Suppose $V$ is finite-dimensional, $T\in\Lin(V,W)$, and $U$ is a subspace of
$W$. Prove that $\{v\in V : Tv\in U\}$ is a subspace of $V$ and
\[ \dim\{v\in V : Tv\in U\}=\dim\nul T+\dim(U\cap\range T). \]
\end{exercise}
\begin{solution}
Let $X=\{v\in V\mid Tv\in U\}$. We first show that $X$ is a subspace of $V$.
Let $v_1,v_2\in X$ and $\alpha,\beta\in\F$. Then
\[ T(\alpha v_1+\beta v_2)=\alpha Tv_1+\beta Tv_2. \]
Since $Tv_1,Tv_2\in U$ and $U$ is a subspace of $W$, it follows that
$\alpha Tv_1+\beta Tv_2\in U$. Therefore, $\alpha v_1+\beta v_2\in X$, and thus
$X$ is a subspace of $V$.

Now, we will compute the dimension of $X$. Let $Y=U\cap\range T$. Let
$\dim T=k$ and $\dim Y=m$. Let $n_1,\dots,n_k$ be a basis of $\nul T$, and let
$y_1,\dots,y_m$ be a basis of $Y$. Since $y_i\in\range T$, there exist vectors
$v_i\in V$ such that $Tv_i=y_i$ for $i=1,\dots,m$. We claim that the set
$\{n_1,\dots,n_k,v_1,\dots,v_m\}$ is a basis of $X$.

First, we show that this set spans $X$. Let $v\in X$. Then $Tv\in U$, and since
$Tv\in\range T$, we have $Tv\in Y$. Therefore, there exist scalars
$\alpha_1,\dots,\alpha_m\in\F$ such that
\[ Tv=\alpha_1y_1+\cdots+\alpha_my_m=\alpha_1Tv_1+\cdots+\alpha_mTv_m. \]
Let $v'=v-(\alpha_1v_1+\cdots+\alpha_mv_m)$. Then
\[ Tv'=Tv-T(\alpha_1v_1+\cdots+\alpha_mv_m)
   =Tv-(\alpha_1Tv_1+\cdots+\alpha_mTv_m)=0. \]
Therefore, $v'\in\nul T$, and we can express $v'$ as a linear combination of the
basis vectors of $\nul T$:
\[ v'=\beta_1n_1+\cdots+\beta_kn_k \]
for some scalars $\beta_1,\dots,\beta_k\in\F$. Thus,
\[ v=v'+(\alpha_1v_1+\cdots+\alpha_mv_m)
   =\beta_1n_1+\cdots+\beta_kn_k+\alpha_1v_1+\cdots+\alpha_mv_m. \]
This shows that $v$ is a linear combination of the vectors in
$\{n_1,\dots,n_k,v_1,\dots,v_m\}$, and hence this set spans $X$.

Next, we show that this set is linearly independent. Suppose
\[ \gamma_1n_1+\cdots+\gamma_kn_k+\delta_1v_1+\cdots+\delta_mv_m=0 \]
for some scalars $\gamma_1,\dots,\gamma_k,\delta_1,\dots,\delta_m\in\F$. Observe
that $Tn_i=0$ for all $i=1,\dots,k$, and $Tv_j=y_j$ for all $j=1,\dots,m$.
Therefore,
\[ 0=T(\gamma_1n_1+\cdots+\gamma_kn_k+\delta_1v_1+\cdots+\delta_mv_m)
   =\delta_1y_1+\cdots+\delta_my_m. \]
Since $y_1,\dots,y_m$ are linearly independent, it follows that
$\delta_1=\cdots=\delta_m=0$. Therefore,
\[ \gamma_1n_1+\cdots+\gamma_kn_k=0. \]
Since $n_1,\dots,n_k$ are linearly independent, $\gamma_1=\cdots=\gamma_k=0$.
Thus, the set $\{n_1,\dots,n_k,v_1,\dots,v_m\}$ is linearly independent and a
basis of $X$. Therefore,
\[ \dim X=k+m=\dim\nul T+\dim(U\cap\range T). \qedhere \]
\end{solution}

\begin{exercise}
Suppose $U$ and $V$ are finite-dimensional vector spaces and $S\in\Lin(V,W)$
and $T\in\Lin(U,V)$. Prove that
\[ \dim\nul ST\le\dim\nul S+\dim\nul T. \]
\end{exercise}
\begin{solution}
Let $u\in U$ such that $u\in\nul ST$. Then $STu=0$, which implies that
$Tu\in\nul S$. Therefore, we can express $\nul ST$ as
\[ \nul ST=\{u\in U\mid Tu\in\nul S\}. \]
Observe that this looks similar to the set in Exercise 3b21, where $T$ is a
linear map from $U$ to $V$ and $\nul S$ is a subspace of $V$. By applying the
result from Exercise 3b21, we have
\[ \dim\nul ST=\dim\nul T+\dim(\nul S\cap\range T). \]
Since $\nul S\cap\range T$ is a subspace of $\nul S$, we have
\[ \dim(\nul S\cap\range T)\le\dim\nul S. \]
Therefore,
\[ \dim\nul ST=\dim\nul T+\dim(\nul S\cap\range T)\le\dim\nul T+\dim\nul S.
   \qedhere \]
\end{solution}

\begin{exercise}
Suppose $U$ and $V$ are finite-dimensional vector spaces and $S\in\Lin(V,W)$
and $T\in\Lin(U,V)$. Prove that
\[ \dim\range ST\le\min\{\dim\range S,\dim\range T\}. \]
\end{exercise}
\begin{solution}
Let $u\in U$ such that $STu\in\range ST$. Then $STu=S(Tu)$, which implies that
$Tu\in\range T$. Therefore, we can express $\range ST$ as
\[ \range ST=\{S(Tu)\mid u\in U\}. \]
Since $Tu\in\range T$, we can rewrite this as
\[ \range ST=\{Sv\mid v\in\range T\}. \]
This shows that $\range ST$ is a subset of $\range S$, and hence
\[ \dim\range ST\le\dim\range S. \]
Now let $v_1,\dots,v_m$ be a basis of $\range T$. Then for any $u\in U$, we can
express $Tu$ as a linear combination of the basis vectors:
\[ Tu=\alpha_1v_1+\cdots+\alpha_mv_m \]
for some scalars $\alpha_1,\dots,\alpha_m\in\F$. Applying $S$ to both sides, we
have
\[ STu=S(Tu)=S(\alpha_1v_1+\cdots+\alpha_mv_m)=\alpha_1Sv_1+\cdots+\alpha_mSv_m. \]
This shows that $\range ST$ is spanned by the vectors $Sv_1,\dots,Sv_m$.
Therefore,
\[ \dim\range ST\le m=\dim\range T. \]
Combining the two inequalities, we have
\[ \dim\range ST\le\min(\dim\range S,\dim\range T). \qedhere \]
\end{solution}

\begin{exercise}\leavevmode
\begin{parts}
\item Suppose $\dim V=5$ and $S,T\in\Lin(V)$ are such that $ST=0$. Prove that
  $\dim\range TS\le2$.
\item Give an example of $S,T\in\Lin(\F^5)$ with $ST=0$ and
  $\dim\range TS=2$.
\end{parts}
\end{exercise}
\begin{solution}\leavevmode
\begin{parts}
\item Since $ST=0$, then $S(Tv)=0$ for all $v\in V$. This implies that
  $\range T\subseteq\nul S$. By the fundamental theorem of linear maps, we have
  \[ 5=\dim\nul S+\dim\range S\ge\dim\range T+\dim\range S. \]
  Now, by Exercise 3B23, we have
  \[ \dim\range TS\le\min(\dim\range T,\dim\range S). \]
  Suppose that $\dim\range TS>2$. Then both $\dim\range T$ and $\dim\range S$
  must be greater than $2$, or equivalently, $\dim\range T\ge3$ and
  $\dim\range S\ge3$. This implies that
  \[ \dim\range T+\dim\range S\ge3+3=6>5, \]
  which is a contradiction. Therefore, $\dim\range TS\le2$.
\item Let $v_1,v_2,v_3,v_4,v_5$ be a basis of $\F^5$. Define
  $S,T\in\Lin(\F^5)$ by
  \[ Sv_1=v_1, \quad Sv_2=v_2, \quad Sv_3=Sv_4=Sv_5=0, \]
  and
  \[ Tv_1=v_3, \quad Tv_2=v_4, \quad Tv_3=Tv_4=Tv_5=0. \]
  Then $ST=0$ since $S(Tv_i)=S(0)=0$ for all $i=1,\dots,5$. Now, we compute
  $\range TS$. We have
  \begin{gather*}
  TSv_1=T(Sv_1)=Tv_1=v_3, \quad TSv_2=T(Sv_2)=Tv_2=v_4,\\
  TSv_3=TSv_4=TSv_5=0.
  \end{gather*}
  Therefore, $\range TS=\spn(v_3,v_4)$, which has dimension $2$. Thus, we have
  constructed an example of $S,T\in\Lin(\F^5)$ with $ST=0$ and
  $\dim\range TS=2$.
\end{parts}
\end{solution}

\begin{exercise}
Suppose that $W$ is finite-dimensional and $S,T\in\Lin(V,W)$. Prove that
$\nul S\subseteq\nul T$ if and only if there exists $E\in\Lin(W)$ such that
$T=ES$.
\end{exercise}
\begin{solution}
\solnote{Manual Remark}{This exercise requires showing that some
$R\in\Lin(\range S,W)$ is well-defined. An issue arises when $T$ is not
injective: for some $w\in\range S$, there may be distinct $v_1,v_2\in V$ such
that $Sv_1=Sv_2=w$. In this case, we must show that $Tv_1=Tv_2$.}

Suppose $\nul S\subseteq\nul T$. Define a linear map $R\colon\range S\to W$ by
$R(Sv)=Tv$ for all $v\in V$. We need to show that $R$ is well-defined. Suppose
$Sv_1=Sv_2$ for some $v_1,v_2\in V$. Then $S(v_1-v_2)=0$, which implies that
$v_1-v_2\in\nul S$. Since $\nul S\subseteq\nul T$, we have
$v_1-v_2\in\nul T$, which means that $T(v_1-v_2)=0$. Therefore, $Tv_1=Tv_2$,
and hence $R(Sv_1)=R(Sv_2)$. Thus, $R$ is well-defined.

By Exercise 3A13, we can extend $R$ to some $E\in\Lin(W)$. Then for any
$v\in V$, we have
\[ Tv=R(Sv)=E(Sv), \]
and hence $T=ES$.

Conversely, suppose there exists $E\in\Lin(W)$ such that $T=ES$. Let
$v\in\nul S$. Then $Sv=0$, and applying $E$ gives
\[ Tv=E(Sv)=E(0)=0. \]
Therefore, $v\in\nul T$, and hence $\nul S\subseteq\nul T$.
\end{solution}

\begin{exercise}
Suppose that $V$ is finite-dimensional and $S,T\in\Lin(V,W)$. Prove that
$\range S\subseteq\range T$ if and only if there exists $E\in\Lin(V)$ such
that $S=TE$.
\end{exercise}
\begin{solution}
\solnote{Manual Remark}{From the previous exercise, the reader may be tempted
to show that a map is well-defined again. However, this is not necessary here,
since we are defining a map on the entire space $V$, not just on $\range S$.}

Suppose $\range S\subseteq\range T$. Let $v_1,\dots,v_n$ be a basis of $V$. For
each $i=1,\dots,n$, since $Sv_i\in\range S\subseteq\range T$, there exists
$u_i\in V$ such that $Tu_i=Sv_i$. Define a linear map $E\in\Lin(V)$ by
\[ Ev_i=u_i \quad\text{for } i=1,\dots,n. \]
Then for any $v\in V$, we can express $v$ as a linear combination of the basis
vectors:
\[ v=\alpha_1v_1+\cdots+\alpha_nv_n \]
for some scalars $\alpha_1,\dots,\alpha_n\in\F$. Applying $S$ and $TE$ to both
sides, we have
\begin{align*}
Sv &= \alpha_1Sv_1+\cdots+\alpha_nSv_n=\alpha_1Tu_1+\cdots+\alpha_nTu_n\\
   &= T(\alpha_1u_1+\cdots+\alpha_nu_n)=TEv.
\end{align*}
Therefore, $S=TE$.

Conversely, suppose there exists $E\in\Lin(V)$ such that $S=TE$. Let
$w\in\range S$. Then there exists $v\in V$ such that $Sv=w$. Applying the
definition of $S$, we have
\[ w=Sv=TE(v)=T\bigl(E(v)\bigr). \]
Therefore, $w\in\range T$, and hence $\range S\subseteq\range T$.
\end{solution}

\begin{exercise}
Suppose $P\in\Lin(V)$ and $P^2=P$. Prove that $V=\nul P\oplus\range P$.
\end{exercise}
\begin{solution}
Let $v\in V$. We can express $v$ as
\[ v=(v-Pv)+Pv. \]
Observe that $Pv\in\range P$ by definition. Now, we need to show that
$v-Pv\in\nul P$. Applying $P$ to $v-Pv$, we have
\[ P(v-Pv)=Pv-P^2v=Pv-Pv=0. \]
Therefore, $v-Pv\in\nul P$. This shows that every vector $v\in V$ can be
expressed as a sum of a vector in $\nul P$ and a vector in $\range P$, which
implies that
\[ V=\nul P+\range P. \]

Next, we need to show that the sum is direct. Suppose
$u\in\nul P\cap\range P$. Then there exists some $w\in V$ such that $u=Pw$.
Since $u\in\nul P$, we have
\[ Pu=P(Pw)=P^2w=Pw=u. \]
But since $u\in\nul P$, we also have $Pu=0$. Therefore, we have
\[ u=Pu=0. \]
This shows that the intersection of $\nul P$ and $\range P$ is trivial, i.e.,
$\nul P\cap\range P=\{0\}$. Hence, the sum is direct, and we conclude that
\[ V=\nul P\oplus\range P. \qedhere \]
\end{solution}

\begin{exercise}
Suppose $D\in\Lin(\Poly(\R))$ is such that $\deg Dp=(\deg p)-1$ for every
nonconstant polynomial $p\in\Poly(\R)$. Prove that $D$ is surjective.
\exnote{The notation $D$ is used above to remind you of the differentiation map
that sends a polynomial $p$ to $p'$.}
\end{exercise}
\begin{solution}
\solnote{Manual Remark}{Note that $\Poly(\R)$ is infinite-dimensional, so we
cannot use the fundamental theorem of linear maps directly. Instead, one should
consider finite-dimensional subspaces of $\Poly(\R)$.}

If $p=0$, then $p=D(0)\in\range D$ by linearity. Suppose $p\ne0$ and let
$n=\deg p$. Consider the list $D(x),D(x^2),\dots,D(x^{n+1})$. Since $x^k$ is
nonconstant for each $k\in\{1,\dots,n+1\}$, we have $\deg D(x^k)=k-1$. Thus,
$D(x),D(x^2),\dots,D(x^{n+1})$ is a list of $n+1$ polynomials with pairwise
distinct degrees, which implies linear independence. Since
$\dim\Poly_n(\R)=n+1$, this list is a basis for $\Poly_n(\R)$. In particular,
$p\in\Poly_n(\R)$, so there exist scalars $c_1,\dots,c_{n+1}\in\R$ such that
\[ p=\sum_{k=1}^{n+1}c_kD(x^k)=D\Biggl(\sum_{k=1}^{n+1}c_kx^k\Biggr)\in\range D. \]
Since $p$ was arbitrary, $D$ is surjective.
\end{solution}

\begin{exercise}
Suppose $p\in\Poly(\R)$. Prove that there exists a polynomial $q\in\Poly(\R)$
such that $5q''+3q'=p$.
\exnote{This exercise can be done without linear algebra, but it's more fun to
do it using linear algebra.}
\end{exercise}
\begin{solution}
\solnote{Appendix}{See the alternate proof at the end of this solution for an
explicit method of determining $q$ given a system of equations.}

Define $D\in\Lin(\Poly(\R))$ by $Dq=5q''+3q'$ for all $q\in\Poly(\R)$. For
nonconstant $q$, we have $\deg Dq=\deg q-1$. By Exercise 3B28, $D$ is
surjective. Therefore, for any polynomial $p\in\Poly(\R)$, there exists a
polynomial $q\in\Poly(\R)$ such that $Dq=p$, which means that $5q''+3q'=p$.

\emph{Alternate proof.} \emph{Coolempire93.} Let
$p(x)=a_0+a_1x+\cdots+a_nx^n$ and suppose there exists a polynomial
$q(x)=b_0+b_1x+\cdots+b_mx^m$ satisfying $5q''+3q'=p$. If $\deg q=m\ge2$, we
can say $q'$ has degree $m-1$ and $q''$ has degree $m-2$, so $5q''+3q'$ has
degree $m-1$. To satisfy the given equation, then, we must have
$m-1=n=\deg p$, and we may write $q(x)=b_0+b_1x+\cdots+b_{n+1}x^{n+1}$. Before
we deal with the general case, we will quickly consider when $n=0$ (which would
force $m<2$):

If $p(x)=a_0$, then let $q(x)=b_0+b_1x$. Because $5q''+3q'=p$, we have
$5q''+3q'=5(0)+3b_1=a_0$, so we set $b_1=a_0/3$ and we have a
polynomial $q$ such that $5q''+3q'=p$ for any value of $b_0$. As we will see
below, the constant term of $q$ never depends on $p$.

For the remainder of this proof, we take $m=n+1$ and suppose $n\ge1$. We have
\begin{align*}
q'(x) &= b_1+2b_2x+\cdots+(n+1)b_{n+1}x^n\\
q''(x) &= 2b_2+6b_3x+\cdots+n(n+1)b_{n+1}x^{n-1}
\end{align*}
Equating coefficients of $5q''+3q'$ with those of $p$, we obtain the system
\[ \begin{cases}
   a_0=3b_1+10b_2\\
   a_1=6b_2+30b_3\\
   \vdots\\
   a_{n-1}=3nb_n+5n(n+1)b_{n+1}\\
   a_n=3(n+1)b_{n+1}
   \end{cases}. \]
That is, $a_i=3(i+1)b_{i+1}+5(i+1)(i+2)b_{i+2}$ for $0\le i<n$ and
$a_n=3(n+1)b_{n+1}$. Finally, we solve for the coefficients of $q$ by
back-substitution. Noting that
\[ b_j=\frac{a_{j-1}-5j(j+1)b_{j+1}}{3j}=\frac{a_{j-1}}{3j}-(5/3)(j+1)b_{j+1}, \]
we have
\[ \begin{cases}
   b_{n+1}=\frac{a_n}{3(n+1)}\\[2pt]
   b_n=\frac{a_{n-1}}{3n}-(5/3)(n+1)b_{n+1}=\frac{a_{n-1}}{3n}-5/9a_n\\[2pt]
   b_{n-1}=\frac{a_{n-2}}{3(n-1)}-5/3nb_n
     =\frac{a_{n-2}}{3(n-1)}-5/9a_{n-1}+25/27na_n\\
   \vdots\\
   b_1=\frac{a_0-10b_2}{3}
   \end{cases} \]
which has a valid solution$^\dagger$ since no division by $0$ occurs. Finally,
we note that $b_0$ is a free variable, so there is an infinite family of
polynomials $q\in\Poly(\R)$ with the given coefficients such that
$5q''+3q'=p$.

$^\dagger$\,Solving this recurrence is outside the scope of this course, but,
specifically, the coefficients of $q$ are
\begin{align*}
b_j &= \frac{a_{j-1}}{3j}-5/9a_j+(25/27)(j+1)a_{j+1}
       -(125/81)(j+1)(j+2)a_{j+2}\\
    &\qquad+\cdots+(-1)^{n-j+1}\frac{5^{n-j+1}}{3^{n-j+2}}(j+1)(j+2)\cdots na_n\\
    &= 1/3\Biggl(\frac{a_{j-1}}{j}
       +\sum_{i=j}^n\Bigl(-5/3\Bigr)^{i-j+1}\frac{i!}{j!}a_i\Biggr).
       \qedhere
\end{align*}
\end{solution}

\begin{exercise}
Suppose $\varphi\in\Lin(V,\F)$ and $\varphi\ne0$. Suppose $u\in V$ is not in
$\nul\varphi$. Prove that
\[ V=\nul\varphi\oplus\{au : a\in\F\}. \]
\end{exercise}
\begin{solution}
Let $v\in V$. We want to show that $v$ can be expressed as a sum of a vector in
$\nul\varphi$ and a scalar multiple of $u$. Since $\varphi(u)\ne0$, we can
define a scalar
\[ a=\frac{\varphi(v)}{\varphi(u)}. \]
Then consider the vector
\[ w=v-au. \]
We have
\[ \varphi(w)=\varphi(v-au)=\varphi(v)-a\varphi(u)
   =\varphi(v)-\frac{\varphi(v)}{\varphi(u)}\varphi(u)=0. \]
Therefore, $w\in\nul\varphi$. Thus, we can express $v$ as
\[ v=w+au, \]
where $w\in\nul\varphi$ and $au$ is a scalar multiple of $u$. This shows that
every vector in $V$ can be expressed as a sum of a vector in $\nul\varphi$ and a
scalar multiple of $u$, which implies that
\[ V=\nul\varphi+\{au\mid a\in\F\}. \]

Next, we need to show that the sum is direct. Suppose
$w\in\nul\varphi\cap\{au\mid a\in\F\}$. Then there exists a scalar $a\in\F$
such that $w=au$. Since $w\in\nul\varphi$, we have
\[ 0=\varphi(w)=\varphi(au)=a\varphi(u). \]
Since $\varphi(u)\ne0$, it follows that $a=0$. Therefore, $w=0$, which shows
that the intersection of $\nul\varphi$ and $\{au\mid a\in\F\}$ is trivial.
Hence, the sum is direct, and we conclude that
\[ V=\nul\varphi\oplus\{au\mid a\in\F\}. \qedhere \]
\end{solution}

\begin{exercise}
Suppose $V$ is finite-dimensional, $X$ is a subspace of $V$, and $Y$ is a
finite-dimensional subspace of $W$. Prove that there exists $T\in\Lin(V,W)$
such that $\nul T=X$ and $\range T=Y$ if and only if $\dim X+\dim Y=\dim V$.
\end{exercise}
\begin{solution}
Suppose there exists $T\in\Lin(V,W)$ such that $\nul T=X$ and $\range T=Y$. By
the fundamental theorem of linear maps, we have
\[ \dim V=\dim\nul T+\dim\range T=\dim X+\dim Y. \]

Conversely, suppose $\dim X+\dim Y=\dim V$. Let $n=\dim V$, $k=\dim X$, and
$m=\dim Y$. Then $n=k+m$. Choose a basis $\{x_1,\dots,x_k\}$ of $X$ and extend
it to a basis $\{x_1,\dots,x_k,v_1,\dots,v_m\}$ of $V$. Let
$\{y_1,\dots,y_m\}$ be a basis of $Y$. Define a linear map $T\in\Lin(V,W)$ by
\[ Tx_i=0 \quad\text{for } i=1,\dots,k, \]
and
\[ Tv_j=y_j \quad\text{for } j=1,\dots,m. \]
Then for any $v\in V$, we can express $v$ as a linear combination of the basis
vectors:
\[ v=\alpha_1x_1+\cdots+\alpha_kx_k+\beta_1v_1+\cdots+\beta_mv_m \]
for some scalars $\alpha_1,\dots,\alpha_k,\beta_1,\dots,\beta_m\in\F$. Applying
$T$ to both sides gives
\begin{align*}
Tv &= \alpha_1Tx_1+\cdots+\alpha_kTx_k+\beta_1Tv_1+\cdots+\beta_mTv_m\\
   &= 0+\beta_1y_1+\cdots+\beta_my_m.
\end{align*}
Therefore, $Tv\in Y$. Now suppose $y=\gamma_1y_1+\cdots+\gamma_my_m\in Y$. Then
\[ T(\gamma_1v_1+\cdots+\gamma_mv_m)=\gamma_1y_1+\cdots+\gamma_my_m=y. \]
Hence, $\range T=Y$. Moreover, if $Tv=0$, then $\beta_1y_1+\cdots+\beta_my_m=0$.
Since $y_1,\dots,y_m$ are linearly independent, it follows that
$\beta_1=\cdots=\beta_m=0$. Thus, $v\in X$. Additionally, if $v\in X$, then
$v=\alpha_1x_1+\cdots+\alpha_kx_k$ for some scalars
$\alpha_1,\dots,\alpha_k\in\F$, and hence $Tv=0$. Therefore, $\nul T=X$. This
shows that there exists a linear map $T\in\Lin(V,W)$ such that $\nul T=X$ and
$\range T=Y$.
\end{solution}

\begin{exercise}
Suppose $V$ is finite-dimensional with $\dim V>1$. Show that if
$\varphi\colon\Lin(V)\to\F$ is a linear map such that
$\varphi(ST)=\varphi(S)\varphi(T)$ for all $S,T\in\Lin(V)$, then $\varphi=0$.
\exnote{Hint: The description of the two-sided ideals of $\Lin(V)$ given by
Exercise 17 in Section 3A might be useful.}
\end{exercise}
\begin{solution}
Observe that $\nul\varphi$ is a two-sided ideal of $\Lin(V)$ as follows: if
$X\in\nul\varphi$ and $Y\in\Lin(V)$, then
\[ \varphi(XY)=\varphi(X)\varphi(Y)=0\cdot\varphi(Y)=0, \quad
   \varphi(YX)=\varphi(Y)\varphi(X)=\varphi(Y)\cdot0=0. \]
Since $\dim V>1$, then Exercise 3A16 gives the existence of $S,T\in\Lin(V)$
such that $ST\ne TS$. Then we have
\[ \varphi(ST-TS)=\varphi(S)\varphi(T)-\varphi(T)\varphi(S)=0. \]
Therefore, $ST-TS\in\nul\varphi$. Since $\nul\varphi$ is a two-sided ideal of
$\Lin(V)$, then by Exercise 3A17, we have $\nul\varphi=\Lin(V)$. Hence,
$\varphi=0$.
\end{solution}

\begin{exercise}
Suppose that $V$ and $W$ are real vector spaces and $T\in\Lin(V,W)$. Define
$T_{\C}\colon V_{\C}\to W_{\C}$ by
\[ T_{\C}(u+iv)=Tu+iTv \]
for all $u,v\in V$.
\begin{parts}
\item Show that $T_{\C}$ is a (complex) linear map from $V_{\C}$ to $W_{\C}$.
\item Show that $T_{\C}$ is injective if and only if $T$ is injective.
\item Show that $\range T_{\C}=W_{\C}$ if and only if $\range T=W$.
\end{parts}
\exnote{See Exercise 8 in Section 1B for the definition of the complexification
$V_{\C}$. The linear map $T_{\C}$ is called the \textbf{complexification} of the
linear map $T$.}
\end{exercise}
\begin{solution}\leavevmode
\begin{parts}
\item Let $u+vi$ and $w+xi\in V_{\C}$. Then for any $a+bi\in\C$, we have
  \begin{align*}
  T_{\C}\bigl((a+bi)(u+vi)+(w+xi)\bigr)
    &= T_{\C}\bigl((au-bv+(av+bu)i)+(w+xi)\bigr)\\
    &= T(au-bv+w)+iT(av+bu+x)\\
    &= aTu-bTv+Tw+i(aTv+bTu+Tx)\\
    &= (a+bi)(Tu+iTv)+(Tw+iTx)\\
    &= (a+bi)T_{\C}(u+vi)+T_{\C}(w+xi).
  \end{align*}
  Therefore, $T_{\C}$ is a complex linear map from $V_{\C}$ to $W_{\C}$.
\item Suppose $T_{\C}$ is injective. Let $v\in V$ be such that $Tv=0$. Then
  \[ T_{\C}(v+i0)=Tv+iT0=0+i0=0. \]
  Since $T_{\C}$ is injective, we have $v+i0=0$, which implies that $v=0$.
  Therefore, $T$ is injective.

  Conversely, suppose $T$ is injective. Let $u+iv\in V_{\C}$ be such that
  $T_{\C}(u+iv)=0$. Then
  \[ Tu+iTv=0. \]
  This implies that $Tu=0$ and $Tv=0$. Since $T$ is injective, we have $u=0$
  and $v=0$. Therefore, $u+iv=0$, and hence $T_{\C}$ is injective.
\item Suppose $\range T_{\C}=W_{\C}$. Let $w\in W$. Then $w+i0\in W_{\C}$, and
  since $\range T_{\C}=W_{\C}$, there exists $u+iv\in V_{\C}$ such that
  \[ T_{\C}(u+iv)=Tu+iTv=w+i0. \]
  This implies that $Tu=w$ and $Tv=0$. Therefore, $w\in\range T$, and hence
  $\range T=W$.

  Conversely, suppose $\range T=W$. Let $w_1+iw_2\in W_{\C}$. Since
  $\range T=W$, there exist $u,v\in V$ such that $Tu=w_1$ and $Tv=w_2$. Then
  \[ T_{\C}(u+iv)=Tu+iTv=w_1+iw_2. \]
  Therefore, $w_1+iw_2\in\range T_{\C}$, and hence $\range T_{\C}=W_{\C}$.
\end{parts}
\end{solution}
